\begin{frame}{Forlægning af vildt}
    \begin{itemize}
        \item \textbf{Betydning}: Partering/udskæring af klovbærende vildt efter opbrækning
        \item Formål: Adskille kød fra knogler, flå skind af, forberede til modning/konsum
        \item Gøres typisk hjemme eller på jagtstedet efter indvolde er fjernet
        \item Grundtrin (eks. rådyr):
          \begin{itemize}
            \item Hæng dyret op (hoved opad) i køligt sted
            \item Flå skindet af (start ved baglår, træk ned mod hoved)
            \item Fjern hoved, ben og skanker i leddene
            \item Skær ud stykker: stege, culotte, fileter, bov, nakke m.m.
            \item Fjern sener, hinder og fedt efter behov
          \end{itemize}
        \item Hygiejne: Rene redskaber, arbejd køligt, undgå forurening
        \item Modning: Hæng kødet 4–10 dage (afhængig af temperatur (ikke over 7 grader), ca. 40 graddøgn)
    \end{itemize}
\end{frame}